\section*{الفصل \ref{ch:g3:food-chains} --- السلاسل الغذائية}

\begin{solution}{exo:g3:food-chains:1}
\omterm{def:g3:food-chains:chain}{``يأكله''}. وهو يشير من المأكول إلى الآكل --- متبعًا الغذاء.
\end{solution}

\begin{solution}{exo:g3:food-chains:2}
\omterm{def:g1:plants-around-us:plant}{نبات} (أو \omterm{def:g1:living-or-not:living}{كائن حي} آخر يصنع غذاءه بنفسه). لأن \omterm{def:g1:plants-around-us:plant}{النباتات} وحدها تتغذى
من غير أن تأكل أحدًا --- من الضوء والماء والهواء؛ وكل حلقة
حيوانية تعيش على تلك الحلقة الخضراء الأولى، مباشرة أو عبر غيرها.
\end{solution}

\begin{solution}{exo:g3:food-chains:3}
عشب $\leftarrow$ جندب $\leftarrow$ سحلية
$\leftarrow$ صقر.
\end{solution}

\begin{solution}{exo:g3:food-chains:4}
الأسهم معكوسة: والصواب عشب $\leftarrow$ أرنب
$\leftarrow$ ثعلب --- فالعشب \omterm{def:g3:food-chains:chain}{يأكله} الأرنب، والأرنب
\omterm{def:g3:food-chains:chain}{يأكله} الثعلب.
\end{solution}

\begin{solution}{exo:g3:food-chains:5}
\omterm{def:g1:plants-around-us:parts}{أوراق} البلوط $\leftarrow$ \omterm{ex:g2:animals-grow-up:butterfly}{يرقة} $\leftarrow$ قرقف أزرق
$\leftarrow$ باشق.
\end{solution}

\begin{solution}{exo:g3:food-chains:6}
الخس \omterm{def:g3:food-chains:chain}{يأكله} الحلزون؛ والقنفذ يأكل الحلزون.
\end{solution}

\begin{solution}{exo:g3:food-chains:7}
مثلًا: \omterm{def:g1:plants-around-us:plant}{نباتات} مائية $\leftarrow$ شرغوف $\leftarrow$
سمكة شوكية $\leftarrow$ بلشون.
\end{solution}

\begin{solution}{exo:g3:food-chains:8}
لأن غذاء غذاء غذائه عشب: فقلة الجنادب على العشب المصفرّ تعني قلة
السحالي، وقلة السحالي تعني صيدًا هزيلًا للصقر. فالتغير في الحلقة
الأولى يسري صعودًا في السلسلة كلها.
\end{solution}

\begin{solution}{exo:g3:food-chains:9}
تسميم الحلزونات يجوّع آكلها: فيرحل القنفذ أو يموت. ثم، وقد ذهب
القنفذ --- عدو الحلزونات الأول --- تتكاثر البزّاقات والحلزونات
الباقية أشد من قبل: فإزالة الحلقات تسري صعودًا في السلسلة (قنفذ
جائع) ونزولًا فيها (حلزونات لا تُصطاد).
\end{solution}

\begin{solution}{exo:g3:food-chains:10}
صحيح --- عبر وسط السلسلة. فالسحلية تأكل الجنادب، والجنادب تأكل
العشب: فإن فشل العشب سرى الفشل صعودًا، حلقةً حلقة، من العشب إلى
الجندب إلى السحلية. والسحلية لا تمسّ العشب قط، وتعتمد عليه مع
ذلك.
\end{solution}
